\documentclass[11pt]{amsart}
\usepackage[margin=1.15in]{geometry}
\usepackage{amsmath,amssymb,amsthm}
\usepackage{mathtools}
\usepackage[hidelinks]{hyperref}

\newtheorem{theorem}{Theorem}[section]
\newtheorem{lemma}[theorem]{Lemma}
\newtheorem{corollary}[theorem]{Corollary}
\newtheorem{proposition}[theorem]{Proposition}
\theoremstyle{remark}
\newtheorem{remark}[theorem]{Remark}

\renewcommand{\thepart}{\Roman{part}}
\newcommand{\R}{\mathbb{R}}
\newcommand{\eps}{\varepsilon}
\DeclareMathOperator{\Hess}{Hess}
\DeclareMathOperator{\arccosh}{arccosh}
\DeclareMathOperator{\arcsinh}{arcsinh}

\title{Maxwell's problem for three point charges}
\author{Ansh Mishra}
\email{ansh.mishra2025@gmail.com}
\author{Aryan Senthilkumar}
\email{aryan.senthilk@gmail.com}
\date{October 24, 2026}

\begin{document}

\begin{abstract}
We study the critical points of the potential $V(x)=\sum_{i=1}^{3}q_i|x-p_i|^{-2\alpha}$ of three point charges.

For charges of one sign we prove that $V$ has at most four critical points, for every $\alpha>0$ and in every dimension
$n\ge2$. The bound is sharp. The best bound known before was six, due to Gabrielov, Novikov, Novikov and Shapiro.

For charges of different signs we prove that the Coulomb potential in $\R^3$ has at most four critical points if the
charges are not collinear, and at most four non-degenerate ones if they are. The same holds for every
$\alpha\in[\tfrac14,2]$. Together, these results prove Maxwell's bound $(k-1)^2$ for $k=3$ Coulomb charges: for all
equilibria when the charges are not collinear, and for the non-degenerate ones when they are.

In the angles under which a critical point sees the sides of the triangle, the equilibrium equations separate.
Whether a critical point is a minimum is decided by an explicit inequality, and two minima are excluded once the region
defined by that inequality is convex. For charges of one sign this convexity follows from log-convexity. For charges of
different signs it fails for large exponents, and we establish it by a rigorous interval-arithmetic computation.
\end{abstract}

\maketitle

\section{Introduction}

In \S113 of his \emph{Treatise} \cite{Max}, Maxwell asserted that $k$ point charges in $\R^3$ produce at most $(k-1)^2$
isolated equilibrium points. His argument is incomplete, and the question became known as Maxwell's problem
\cite{GNS}.

Gabrielov, Novikov and Shapiro \cite{GNS} proved the first finite bounds, using fewnomial theory. In particular, they
showed that three positive charges have at most $12$ non-degenerate equilibria, for every exponent $\alpha>0$.

Very recently:
\begin{itemize}
\item Arathoon, Ball and Kvalheim \cite{ABK} constructed five positive charges with at least $24>16$ non-degenerate
  equilibria. So Maxwell's bound fails in general.
\item Gabrielov, Novikov, Novikov and Shapiro \cite{GNNS} improved $12$ to $6$ for three positive charges. They note
  that Maxwell's value $4$ was known only for three charges of equal magnitudes \cite{Tsai}.
\end{itemize}
See also \cite{EFO,EP,Oli,Zol} for related work. We settle the three-charge case for charges of one sign and every
exponent. For charges of different signs and the Coulomb potential we prove the same bound for all equilibria of
non-collinear charges, and for the non-degenerate equilibria of collinear ones (Corollary~\ref{cor:maxwell}).

\begin{theorem}[charges of one sign]\label{thm:one}
Let $\alpha>0$, $n\ge 2$, $q_1,q_2,q_3>0$, and let $p_1,p_2,p_3\in\R^n$ be distinct. Then the potential
\[
  V(x)=\sum_{i=1}^{3}q_i\,|x-p_i|^{-2\alpha}
\]
has at most four critical points in $\R^n\setminus\{p_1,p_2,p_3\}$.
\begin{itemize}
\item If the $p_i$ are collinear, there are exactly two critical points.
\item Otherwise, every critical point lies in the open triangle $T=p_1p_2p_3$. At most one of them is a local minimum of
  $V|_\Pi$ or degenerate, where $\Pi$ is the plane of the charges. All the others are non-degenerate saddles of
  $V|_\Pi$, and there are at most three of them.
\end{itemize}
\end{theorem}

The bound is attained, e.g.\ by three equal charges at the vertices of an equilateral triangle: the centre is a minimum
in $\Pi$, and there are three saddles. Coulomb's law in $\R^3$ is the case $\alpha=1/2$.

\begin{theorem}[charges of different signs]\label{thm:mixed}
Let $q_1,q_2,q_3$ be nonzero real numbers, not all of the same sign, placed at distinct points $p_1,p_2,p_3\in\R^3$. Let
$V(x)=\sum_iq_i|x-p_i|^{-1}$.
\begin{enumerate}
\item[(a)] If the $p_i$ are not collinear, then $V$ has at most four critical points, and all of them lie in the plane of
  the charges.
\item[(b)] If the $p_i$ are collinear, then the critical points off the line form circles about it, and $V$ has at most
  four non-degenerate critical points.
\end{enumerate}
\end{theorem}

\begin{corollary}[Maxwell's bound for three charges]\label{cor:maxwell}
Three point charges in $\R^3$ with the Coulomb potential have at most four non-degenerate equilibrium points. If the
charges are not collinear, they have at most four equilibrium points altogether.
\end{corollary}

\begin{proof}
For charges of one sign this is Theorem~\ref{thm:one} with $\alpha=\tfrac12$, which counts all critical points.
Otherwise it is Theorem~\ref{thm:mixed}.
\end{proof}

\begin{theorem}[other exponents]\label{thm:alpha}
Parts (a) and (b) of Theorem~\ref{thm:mixed} also hold for $V=\sum_iq_i|x-p_i|^{-2\alpha}$ for every
$\alpha\in[\tfrac14,2]$: part (a) in $\R^n$ for every $n\ge2$, and part (b) in $\R^n$ for $n\ge3$, where the critical
points off the line form $(n-2)$-spheres about it.
\end{theorem}
For $\alpha=1,\tfrac32,2$ these are the Newtonian potentials of $\R^4,\R^5,\R^6$.

The proof of Theorem~\ref{thm:one} does not use a computer. Theorems~\ref{thm:mixed} and~\ref{thm:alpha} rely on one
computer-assisted step, Proposition~\ref{prop:M}, which is verified by interval arithmetic.

\medskip
\noindent\textbf{Idea of the proof.}
For a point $x$ there are unique charges, up to a common factor, for which $x$ is an equilibrium. The proof makes three
observations about them.
\begin{enumerate}
\item They take a separated form in the angles under which $x$ sees the sides of the triangle; see \eqref{eq:sep}. This
  rests on Euler's theorem on pedal triangles. It turns the equilibrium equations into the gradient equations of a
  function whose terms each depend on one angle.
\item Whether $x$ is a minimum is decided by an explicit inequality $E<(b+1)^2/b$ in these angles, with $b=2\alpha+1$
  (Lemma~\ref{lem:D}).
\item Where this inequality holds, the separated function is strictly convex (or concave). So two minima for the same
  charges are impossible once the region it defines is convex.
\end{enumerate}
For charges of one sign, $E$ is a sum of products of log-convex functions, so the region is convex
(Lemma~\ref{lem:convex}), and a Morse count gives at most three saddles.

For charges of different signs, the minima lie outside the triangle: beyond the side opposite the odd charge, outside
the circumcircle, around the excentre. One product now enters $E$ with a minus sign, and the region is not convex for
large exponents (Remark~\ref{rem:large}). We describe its boundary in closed form through a strip model of the
hyperbolic plane. We then reduce its convexity to the positivity of an explicit function of three variables (four, when
the exponent varies), and verify that positivity by interval arithmetic (Section~\ref{sec:cert}).

Part~I (Sections~\ref{sec:red}--\ref{sec:rem1}) proves Theorem~\ref{thm:one}. Part~II
(Sections~\ref{sec:setting}--\ref{sec:rem2}) proves Theorems~\ref{thm:mixed} and~\ref{thm:alpha}.

\part{Charges of one sign}

\section{Reduction}\label{sec:red}

Throughout Part~I, $\beta=2\alpha+2$, $b=\beta-1=2\alpha+1>1$, and
\[
  c=\frac{\beta-1}{\beta^2}=\frac{b}{(b+1)^2}.
\]
Write $e_2(v)=v_1v_2+v_2v_3+v_3v_1$ for $v\in\R^3$.

\begin{lemma}[location]\label{lem:loc}
Every critical point of $V$ is of the form $x=\sum w_ip_i/\sum w_i$ with $w_i=q_i|x-p_i|^{-\beta}>0$.
\begin{itemize}
\item If the $p_i$ are collinear, $V$ has exactly two critical points.
\item Otherwise, the critical points of $V$ are exactly the critical points of $V|_\Pi$ that lie in the open triangle
  $T$.
\end{itemize}
\end{lemma}

\begin{proof}
We have $\nabla V(x)=-2\alpha\sum_i q_i r_i^{-\beta}(x-p_i)$, with $r_i=|x-p_i|$. This gives the first claim.

\emph{Collinear case.} Every critical point lies on the open segment spanned by the charges, and not at a charge. On each
of the two open intervals between consecutive charges, $V$ is strictly convex (each $|t-s|^{-2\alpha}$ is) and tends to
$+\infty$ at both ends. So each interval contains exactly one critical point.

\emph{Non-collinear case.} The point $x$ is a combination of the $p_i$ with positive weights, so it lies in $T$. For
$x\in\Pi$ the gradient lies in $\Pi$, so critical points of $V$ in $\Pi$ are the same as critical points of $V|_\Pi$.
\end{proof}

From now on the charges are not collinear, and we work in $\Pi=\R^2$. For $x\in T$ let:
\begin{itemize}
\item $\lambda_i=\lambda_i(x)$ be the barycentric coordinates, $r_i=|x-p_i|$, and $u_i=(x-p_i)/r_i$;
\item $\theta_i=\angle p_jxp_k$ be the angle at $x$ subtending the side opposite $p_i$, where $\{i,j,k\}=\{1,2,3\}$;
\item $\Delta$ be the area of $T$, $A_i$ its angle at $p_i$, $a_i=|p_jp_k|$, and $R$, $O$ its circumradius and
  circumcentre.
\end{itemize}
Then $\theta_1+\theta_2+\theta_3=2\pi$ and $\theta_i<\pi$. Since $\Delta\lambda_i$ is the area of $xp_jp_k$,
\begin{equation}\label{eq:lambda}
  \lambda_i=\frac{r_jr_k\sin\theta_i}{2\Delta}.
\end{equation}
Put
\[
  S(x)=\sum_i\lambda_iu_iu_i^{\mathsf T}\quad(\operatorname{tr}S=1).
\]
By the Cauchy--Binet formula,
\begin{equation}\label{eq:D}
  D(x):=\det S(x)=\sum_{i<j}\lambda_i\lambda_j\,(u_i\times u_j)^2
        =\sum_{\{i,j,k\}}\lambda_i\lambda_j\sin^2\theta_k .
\end{equation}

\begin{lemma}[reduction]\label{lem:red}
Let $x\in T$. Then $x$ is a critical point of $V=V_q$ if and only if $q$ is proportional to
$Q(x):=(\lambda_ir_i^\beta)_{i=1}^3$. In that case $x$ is
\begin{itemize}
\item a non-degenerate local minimum of $V|_\Pi$ if $D(x)>c$;
\item a non-degenerate saddle if $D(x)<c$;
\item degenerate if $D(x)=c$.
\end{itemize}
$V|_\Pi$ has no local maxima.
\end{lemma}

\begin{proof}
By Lemma~\ref{lem:loc} and uniqueness of barycentric coordinates, $x$ is critical if and only if
$(q_ir_i^{-\beta})_i$ is proportional to $(\lambda_i)_i$.

In that case write $q_ir_i^{-\beta}=\kappa\lambda_i$ with $\kappa>0$. Then
\[
  \Hess V(x)=-2\alpha\sum_iq_ir_i^{-\beta}\bigl(I-\beta u_iu_i^{\mathsf T}\bigr)=2\alpha\kappa(\beta S-I).
\]
If $s\le 1/2$ is the smaller eigenvalue of $S$, the eigenvalues of $\beta S-I$ are $\beta s-1$ and $\beta(1-s)-1>0$,
using $\beta>2$. The function $t\mapsto t(1-t)$ is increasing on $[0,1/2]$. Hence $\beta s-1>0$ holds if and only if
\[
  D=s(1-s)>\tfrac1\beta\bigl(1-\tfrac1\beta\bigr)=c.
\]
Likewise $\beta s-1<0$ if and only if $D<c$, and $\beta s-1=0$ if and only if $D=c$.
\end{proof}

\begin{lemma}[index count]\label{lem:index}
If all critical points of $V|_\Pi$ are non-degenerate, then
\[
  \#\{\text{saddles}\}=\#\{\text{minima}\}+2.
\]
\end{lemma}

\begin{proof}
Let $\Omega$ be a large disc minus small discs around the $p_i$.
\begin{itemize}
\item On the outer circle, $\nabla V=-2\alpha(\sum q_i)|x|^{-\beta}x\,(1+o(1))$ has degree $1$.
\item On the circle around $p_i$, $\nabla V=-2\alpha q_ir_i^{-\beta}(x-p_i)(1+o(1))$ also has degree $1$.
\end{itemize}
By the Poincar\'e--Hopf theorem, the indices of the zeros of $\nabla V$ in $\Omega$ sum to $1-3=-2$. Minima have
index $+1$, saddles $-1$, and there are no maxima by Lemma~\ref{lem:red}.
\end{proof}

\section{Angle coordinates and separation of variables}

\begin{lemma}\label{lem:inj}
The map $x\mapsto\theta(x)=(\theta_1,\theta_2,\theta_3)$ is injective on $T$.
\end{lemma}

\begin{proof}
By the inscribed angle theorem, the points $y\in T$ with $\angle p_2yp_3=\theta_1$ lie on a circle $\Gamma_1$ through
$p_2$ and $p_3$. Similarly, those with $\angle p_3yp_1=\theta_2$ lie on a circle $\Gamma_2$ through $p_3$ and $p_1$.

If $\Gamma_1=\Gamma_2$, this circle passes through $p_1,p_2,p_3$, so it is the circumcircle, which does not meet the
open triangle. Hence $\Gamma_1\ne\Gamma_2$. Two distinct circles meet in at most two points, and one of them is
$p_3\notin T$.
\end{proof}

\begin{lemma}[pedal identity]\label{lem:pedal}
Let $x\in T$, put $\delta_i=\theta_i-A_i$, and let $P_x=R^2-|x-O|^2>0$. Then
\begin{enumerate}
\item[(a)] $\delta_i>0$ and $\delta_1+\delta_2+\delta_3=\pi$;
\item[(b)] $r_jr_k\sin\delta_i=P_x\sin A_i$ for $\{i,j,k\}=\{1,2,3\}$.
\end{enumerate}
\end{lemma}

The $\delta_i$ are the angles of the pedal triangle of $x$. Part~(b) is equivalent to Euler's formula
$\Delta\,P_x/(4R^2)$ for the area of the pedal triangle \cite{Johnson}. For completeness we give a short proof.

\begin{proof}
Let $\iota$ be the inversion in the unit circle centred at $x$, and put $p_i^*=\iota(p_i)$. We use three standard
properties of $\iota$:
\begin{enumerate}
\item[(i)] $|\iota(y)-\iota(z)|=|y-z|/(|y-x|\,|z-x|)$;
\item[(ii)] the triangles $xyz$ and $x\,\iota(z)\,\iota(y)$ are similar;
\item[(iii)] a circle with centre $O$ and radius $R$ not passing through $x$ is mapped to a circle of radius
  $R/|R^2-|x-O|^2|$.
\end{enumerate}

Since $\iota$ preserves the rays from $x$, we have $\angle p_j^*xp_k^*=\theta_i$. These three angles are less than $\pi$
and sum to $2\pi$, so $x$ lies inside $T^*=p_1^*p_2^*p_3^*$. By (ii), the angle of $T^*$ at $p_i^*$ is
\[
  \angle p_j^*p_i^*x+\angle xp_i^*p_k^*=\angle xp_jp_i+\angle xp_kp_i
  =(A_j-\angle xp_jp_k)+(A_k-\angle xp_kp_j)=A_j+A_k-(\pi-\theta_i)=\delta_i .
\]
This proves (a).

The circumcircle of $T$ does not pass through $x$, since $x$ lies inside it. By (iii) it is mapped to the circumcircle
of $T^*$, of radius $R/P_x$. The law of sines in $T^*$ and in $T$, together with (i), gives
\[
  \frac{a_i}{r_jr_k}=|p_j^*p_k^*|=\frac{2R}{P_x}\sin\delta_i,
  \qquad a_i=2R\sin A_i .
\]
This is (b).
\end{proof}

\begin{corollary}[separated form]\label{cor:sep}
For $x\in T$,
\begin{equation}\label{eq:sep}
  \lambda_ir_i^\beta=C(x)\,\frac{\sin\theta_i\,\sin^b(\theta_i-A_i)}{\sin^bA_i},
  \qquad C(x)=\frac{(r_1r_2r_3)^{b+1}}{2\Delta\,P_x^{\,b}}>0 .
\end{equation}
\end{corollary}

\begin{proof}
By Lemma~\ref{lem:pedal}(b),
\[
  r_i=\frac{r_1r_2r_3}{r_jr_k}=\frac{r_1r_2r_3\sin\delta_i}{P_x\sin A_i},
\]
and by \eqref{eq:lambda}, $\lambda_ir_i=r_1r_2r_3\sin\theta_i/(2\Delta)$. Now write
$\lambda_ir_i^\beta=\lambda_ir_i\cdot r_i^{\,b}$.
\end{proof}

\section{The region of minima}

Let $L=\{\theta\in\R^3:\theta_1+\theta_2+\theta_3=2\pi\}$ and let
\[
  H_A=\{\theta\in L:\ A_i<\theta_i<\pi\ (i=1,2,3)\},
\]
an open convex polygon. By Lemma~\ref{lem:pedal}(a), $\theta(T)\subset H_A$. For $\theta\in H_A$ define
\[
  W_i(\theta_i)=\frac{\sin A_i}{\sin(\theta_i-A_i)\sin\theta_i}=\cot(\theta_i-A_i)-\cot\theta_i>0,
  \qquad E(\theta)=e_2(W_1,W_2,W_3).
\]

\begin{lemma}\label{lem:D}
For every $x\in T$,
\[
  D(x)=\frac{P_x\sin\theta_1\sin\theta_2\sin\theta_3}{2\Delta}=\frac{1}{E(\theta(x))}.
\]
Consequently $D(x)>c$ if and only if $E(\theta(x))<(b+1)^2/b$.
\end{lemma}

\begin{proof}
By Lemma~\ref{lem:pedal}(b), $W_i=r_jr_k/(P_x\sin\theta_i)$. Hence $W_iW_j=r_1r_2r_3\,r_k/(P_x^2\sin\theta_i\sin\theta_j)$ and
\[
  E=\frac{r_1r_2r_3}{P_x^2\,\prod_l\sin\theta_l}\sum_kr_k\sin\theta_k .
\]
By \eqref{eq:lambda} and \eqref{eq:D},
\[
  D=\frac{r_1r_2r_3\prod_l\sin\theta_l}{(2\Delta)^2}\sum_kr_k\sin\theta_k .
\]
Therefore
\[
  DE=\Bigl(\frac{r_1r_2r_3\sum_kr_k\sin\theta_k}{2\Delta\,P_x}\Bigr)^2 .
\]
Since $r_ir_j\sin\theta_k=2\Delta\lambda_k$, we get $r_1r_2r_3\sum_kr_k\sin\theta_k=2\Delta\sum_k\lambda_kr_k^2$. Finally
\[
  \sum_k\lambda_k|x-p_k|^2=|x-O|^2-2\Bigl\langle x-O,\sum_k\lambda_k(p_k-O)\Bigr\rangle+R^2=R^2-|x-O|^2=P_x .
\]
So $r_1r_2r_3\sum_kr_k\sin\theta_k=2\Delta P_x$. Substituting this into the expressions for $D$ and $E$ gives
$D=P_x\prod_l\sin\theta_l/(2\Delta)$ and $DE=1$.
\end{proof}

\begin{lemma}[convexity]\label{lem:convex}
$E$ is strictly convex on $H_A$, i.e.\ along every non-degenerate segment in $H_A$. Consequently:
\begin{itemize}
\item $\widehat M=\{\theta\in H_A:E(\theta)<(b+1)^2/b\}$ is convex;
\item if $E(\theta),E(\theta')\le(b+1)^2/b$, then the open segment between $\theta$ and $\theta'$ lies in $\widehat M$.
\end{itemize}
\end{lemma}

\begin{proof}
Put $\ell_i(t)=\log W_i(t)=\log\sin A_i-\log\sin(t-A_i)-\log\sin t$. Then
\[
  \ell_i''(t)=\csc^2(t-A_i)+\csc^2t>0 .
\]
For $i\ne j$, $W_iW_j=\exp(\ell_i(\theta_i)+\ell_j(\theta_j))$ is convex, and its Hessian is at least
$W_iW_j\Hess(\ell_i+\ell_j)$.

A nonzero direction $v$ with $\sum v_i=0$ has two nonzero entries, say $v_i$ and $v_j$. Then
$\ell_i''v_i^2+\ell_j''v_j^2>0$, so $E$ is strictly convex along $v$. The two consequences follow.
\end{proof}

\begin{lemma}[definiteness]\label{lem:def}
For $\theta\in H_A$ let $d_i=\cot\theta_i+b\cot(\theta_i-A_i)$. Then
\[
  e_2(d)=(b+1)^2-b\,E(\theta).
\]
If $\theta\in\widehat M$, then $\sum_id_iv_i^2>0$ for all nonzero $v\in\R^3$ with $\sum v_i=0$.
\end{lemma}

\begin{proof}
If three angles in $(0,\pi)$ sum to $\pi$ or to $2\pi$, their cotangents $C$ satisfy $e_2(C)=1$. Indeed,
$\cot(\theta_1+\theta_2)=-\cot\theta_3$ unwinds to this.

Let $X=(\cot\delta_i)_i$ and $C=(\cot\theta_i)_i$, so $e_2(X)=e_2(C)=1$ by Lemma~\ref{lem:pedal}(a). Put
$B=\sum_{i\ne j}X_iC_j$. Since $W=X-C$ and $d=C+bX$,
\[
  e_2(W)=2-B,\qquad e_2(d)=1+b^2+bB=(b+1)^2-b\,e_2(W).
\]

In the basis $(e_1-e_3,e_2-e_3)$ of $\{\sum v_i=0\}$, the form $\sum d_iv_i^2$ has Gram matrix
\[
  G=\begin{pmatrix}d_1+d_3&d_3\\ d_3&d_2+d_3\end{pmatrix},\qquad \det G=e_2(d)>0\ \text{ on }\widehat M .
\]
So $G$ is definite on $\widehat M$. Since $\widehat M$ is connected (Lemma~\ref{lem:convex}), the eigenvalues of $G$
never vanish there and have constant signs.

At $\theta^\circ=((\pi+A_i)/2)_i$, which is the $\theta$-value of the incentre, we have
$W_i=2\tan(A_i/2)$. Hence
\[
  E(\theta^\circ)=4\sum_{i<j}\tan\tfrac{A_i}2\tan\tfrac{A_j}2=4<\frac{(b+1)^2}{b},
\]
so $\theta^\circ\in\widehat M$. Moreover $d_i=(b-1)\tan(A_i/2)>0$ there, so $G$ is positive definite on $\widehat M$.
\end{proof}

\section{Proof of Theorem~\ref{thm:one}}\label{sec:proof1}

The collinear case is Lemma~\ref{lem:loc}. Let the charges be non-collinear and fix $q$. On $H_A$ put
\[
  f_i(t)=\log\sin t+b\log\sin(t-A_i),\qquad f(\theta)=(f_1(\theta_1),f_2(\theta_2),f_3(\theta_3)),
\]
and let $\psi_q=(\log q_i+b\log\sin A_i)_i$. Note that $f_i'=d_i$.

\medskip\noindent\textbf{Step 1: dictionary.} Let $x\in T$ be a critical point of $V_q$ and $\theta=\theta(x)$.
\begin{itemize}
\item By Lemma~\ref{lem:red} and \eqref{eq:sep}, $f(\theta)-\psi_q\in\R(1,1,1)$.
\item By Lemmas~\ref{lem:red} and~\ref{lem:D}, $x$ is a minimum of $V_q|_\Pi$ or degenerate if and only if
  $D(x)\ge c$, if and only if $E(\theta)\le(b+1)^2/b$.
\end{itemize}

\medskip\noindent\textbf{Step 2: at most one critical point with $D\ge c$.} Suppose $x\ne x'$ are two such points, with
$\theta=\theta(x)$ and $\theta'=\theta(x')$.
\begin{itemize}
\item By Lemma~\ref{lem:inj}, $v=\theta-\theta'\ne0$, and $\sum v_i=0$.
\item By Step~1, $f(\theta)-f(\theta')=\mu(1,1,1)$ for some $\mu\in\R$.
\end{itemize}
Hence
\[
  0=\mu\sum_iv_i=\sum_i\bigl(f_i(\theta_i)-f_i(\theta_i')\bigr)v_i=\int_0^1\sum_id_i(\theta_i'+tv_i)\,v_i^2\,dt .
\]
For $0<t<1$ the point $\theta'+tv$ lies in $\widehat M$ (Lemma~\ref{lem:convex}), so the integrand is positive
(Lemma~\ref{lem:def}). This is a contradiction.

\medskip\noindent\textbf{Step 3: at most three saddles.}
Let $\Psi=(\log(Q_1/Q_3),\log(Q_2/Q_3))\colon T\to\R^2$.

\emph{(a) Degeneracy is singularity of $D\Psi$.} Let $x_0$ be a critical point of $V_q$. Put $\eps_3=0$ and, for
$i=1,2$,
\[
  \eps_i(x)=\log(q_i/q_3)-\log(Q_i(x)/Q_3(x)) .
\]
Then $q_i=\kappa(x)Q_i(x)e^{\eps_i(x)}$ with $\kappa=q_3/Q_3>0$. Using $Q_ir_i^{-\beta}=\lambda_i$ and
$\sum\lambda_i(x-p_i)=0$,
\[
  \nabla V_q(x)=-2\alpha\kappa(x)\sum_i\lambda_i(x)\bigl(e^{\eps_i(x)}-1\bigr)(x-p_i).
\]
At $x_0$ we have $\eps=0$, so
\[
  \Hess V_q(x_0)=2\alpha\kappa\,\mathcal L\circ D\Psi(x_0),\qquad
  \mathcal L(e_1,e_2)=\lambda_1e_1(x_0-p_1)+\lambda_2e_2(x_0-p_2).
\]
The map $\mathcal L$ is invertible because $x_0,p_1,p_2$ are not collinear. So $x_0$ is degenerate if and only if
$D\Psi(x_0)$ is singular.

\emph{(b) Density of Morse charges.} By Sard's theorem the regular values of $\Psi$ are dense in $\R^2$. Charges whose
log-ratios form a regular value give a Morse function $V|_\Pi$.

\emph{(c) Persistence of saddles.} Suppose $V_q$ had four distinct saddles $x_1,\dots,x_4$. They are non-degenerate and
satisfy $D<c$. By the implicit function theorem applied to $\Psi$ at the $x_j$, every charge vector $q'$ close enough
to $q$ has critical points $x_1',\dots,x_4'$ near $x_1,\dots,x_4$. These are pairwise distinct and satisfy $D(x_j')<c$,
so they are four saddles of $V_{q'}$.

Choose such a $q'$ for which $V_{q'}|_\Pi$ is Morse. By Step~2, $V_{q'}$ has at most one minimum, so by
Lemma~\ref{lem:index} it has at most three saddles. This is a contradiction, so $V_q$ has at most three saddles.

\medskip\noindent\textbf{Conclusion.} By Lemma~\ref{lem:red}, every critical point of $V_q|_\Pi$ either has $D\ge c$ or is a
saddle. By Steps~2 and~3 there are at most $1+3=4$ critical points. By Lemma~\ref{lem:loc} these are all the critical
points of $V$ in $\R^n\setminus\{p_1,p_2,p_3\}$. \qed

\section{Remarks on Part~I}\label{sec:rem1}

\begin{remark}[a hyperbolic reading]\label{rem:hyp}
The cotangent vectors of the angles of a triangle lie on the sheet $\{e_2=1,\ \sum>0\}$. Equip $\R^3$ with the
polarisation $\langle u,v\rangle=\frac12\sum_{i\ne j}u_iv_j$ of $e_2$. Then this sheet is a hyperboloid model of the
hyperbolic plane. Indeed, the triangle with vertices $0,1,z$ ($\operatorname{Im}z>0$) has cotangent vector
\[
  \bigl(\operatorname{Re}z,\ 1-\operatorname{Re}z,\ |z|^2-\operatorname{Re}z\bigr)/\operatorname{Im}z,
\]
and this map is an isometry from the Poincar\'e half-plane.

Take $X=(\cot\delta_i)$, the shape of the pedal triangle, and $Y=(\cot(\pi-\theta_i))$, the shape of the triangle with
angles $\pi-\theta_i$. The proof of Lemma~\ref{lem:def} shows that $E=2+2\langle X,Y\rangle=2+2\cosh\rho$, where $\rho$ is
the hyperbolic distance between these two shapes. Thus
\[
  D=\frac{1}{4\cosh^2(\rho/2)},
\]
and $x$ is a minimum for its charges if and only if $\rho<\log(2\alpha+1)$. The classical quantity
$|Z|$, $Z=\sum_i\lambda_ie^{2i\arg u_i}$, $D=(1-|Z|^2)/4$, equals $\tanh(\rho/2)$.
\end{remark}

\begin{remark}
(i) \emph{The angle map is a bijection $\theta\colon T\to H_A$, with explicit inverse}
\[
  \lambda_i\propto\frac{\sin A_i\sin\theta_i}{\sin(\theta_i-A_i)} .
\]
Indeed, by \eqref{eq:lambda} and Lemma~\ref{lem:pedal}(b), every $x\in T$ satisfies
\[
  \lambda_i=\frac{P_x\sin A_i\sin\theta_i}{2\Delta\sin\delta_i}.
\]
Given $\theta\in H_A$, the weights $\sin A_i\sin\theta_i/\sin(\theta_i-A_i)$ are positive, so they define a point
$x\in T$. Then $\theta(x)=\theta$, because $\theta\mapsto[\sin\theta_i/\sin(\theta_i-A_i)]_i$ is injective on $H_A$:
each coordinate is decreasing in $\theta_i$, and $\sum\theta_i$ is fixed.

(ii) \emph{The incentre is the unique maximum point of $D$, where $D=1/4$.} Each $W_i$ has a critical point at
$\theta_i^\circ=(\pi+A_i)/2$, so $\theta^\circ$ is a critical point of the strictly convex function $E$ on $H_A$. Hence
it is the unique minimum point of $E$, with $E(\theta^\circ)=4$.

(iii) \emph{$V_q$ has a non-degenerate planar minimum if and only if $\psi_q\in f(\widehat M)+\R(1,1,1)$.} This follows
from Step~1 and (i). In that case all other critical points have $D<c$ by Step~2. So $V_q|_\Pi$ is Morse, and by
Lemma~\ref{lem:index} it has exactly four critical points. Moreover $f$ is injective on $\widehat M$ by Step~2.
\end{remark}

\begin{remark}
We checked all identities used in Part~I numerically at random triangles and points, with relative errors below
$10^{-9}$. We also compared the predictions of Theorem~\ref{thm:one} and Lemma~\ref{lem:D} with direct critical-point
searches, over several thousand configurations and several values of $\alpha$, and found no discrepancy.
\end{remark}

\part{Charges of different signs}

\section{Setting}\label{sec:setting}

In Part~II, write $\alpha=\tfrac12$ (or $\alpha\in[\tfrac14,2]$ for Theorem~\ref{thm:alpha}), $b=2\alpha+1$, $\beta=b+1$,
$c=b/(b+1)^2$ and $K=(b+1)^2/b$. Put $F=K/2-1$; for Coulomb, $K=\tfrac92$ and $F=\tfrac54$.

Replacing $V$ by $-V$ and relabelling, we may assume $q_1,q_2>0>q_3$. The triangle $p_1p_2p_3$ is labelled
counter-clockwise. It has angles $A_i$, area $\Delta$, circumradius $R$ and circumcentre $O$. Put
\[
  P_x=R^2-|x-O|^2 .
\]
For $x$ off the side lines, let $\lambda_i$ be the signed barycentric coordinates and $r_i=|x-p_i|$. Let
$\sigma_i\in(-\pi,\pi)$ be the signed angle at $x$ from $p_j-x$ to $p_k-x$, for $(i,j,k)$ cyclic, so that
$\lambda_i=r_jr_k\sin\sigma_i/(2\Delta)$. Write $e_2(v)=v_1v_2+v_2v_3+v_3v_1$.

\begin{lemma}[planarity and type]\label{lem:type}
Let the charges be non-collinear.
\begin{enumerate}
\item Every critical point of $V$ lies in the plane $\Pi$ of the charges.
\item A point $x\in\Pi$ is critical if and only if $(q_ir_i^{-\beta})_i=\kappa(\lambda_i)_i$ for some $\kappa\ne0$.
\item Write $S=\sum_i\lambda_iu_iu_i^{\mathsf T}$ (so $\operatorname{tr}S=1$) with $u_i=(x-p_i)/r_i$, and $D=\det S$. Then
  $\Hess V|_\Pi(x)=2\alpha\kappa(\beta S-I)$. Hence $x$ is an extremum of $V|_\Pi$ if and only if $D>c$: a minimum if
  $\kappa>0$, a maximum if $\kappa<0$. If $D<c$, then $x$ is a non-degenerate saddle.
\end{enumerate}
\end{lemma}

\begin{proof}
(1) Write $x=y+hn$ with $y\in\Pi$ and $n\perp\Pi$, and put $w_i=q_ir_i^{-\beta}\ne0$. Then
$\nabla V=-2\alpha\bigl(\sum w_i(y-p_i)+hn\sum w_i\bigr)$. If $h\ne0$, then $\sum w_i=0$ and $\sum w_ip_i=0$, which forces
$w=0$ because the $p_i$ are affinely independent; this is impossible.

(2) This is the same computation, by uniqueness of barycentric coordinates.

(3) $\Hess V=-2\alpha\sum w_i(I-\beta u_iu_i^{\mathsf T})$. The eigenvalues of $S$ are $s$ and $1-s$; if $s\le\tfrac12$,
then $\beta(1-s)-1>0$ because $\beta>2$. So the Hessian is definite if and only if $s>1/\beta$, which is equivalent to
$D=s(1-s)>c$.
\end{proof}

\begin{lemma}[signed identities]\label{lem:signed}
For every $x$ off the side lines:
\begin{enumerate}
\item[(P)] $r_jr_k\sin(\sigma_i-A_i)=P_x\sin A_i$;
\item[(D)] $D(x)=P_x\sin\sigma_1\sin\sigma_2\sin\sigma_3/(2\Delta)$;
\item[(E)] $D(x)\,e_2(W)=1$, where $W_i=\cot(\sigma_i-A_i)-\cot\sigma_i$.
\end{enumerate}
\end{lemma}

\begin{proof}
With $p_1=(0,0)$, $p_2=(1,0)$ and $p_3=(u,v)$, each identity becomes a polynomial identity in $x$, $u$, $v$ and the
square roots $r_i$ and $|p_j-p_k|$, once these roots are treated as symbols subject to their squares.
\begin{itemize}
\item For (P), multiply by $|p_i-p_j||p_i-p_k|$.
\item For (D), multiply by $(r_1r_2r_3)^2$.
\item For (E), clear denominators.
\end{itemize}
Expanding, each difference vanishes identically (\texttt{code/symbolic\_mixed.py}). For $x$ inside the triangle these
are Lemma~\ref{lem:pedal}(b) and Lemma~\ref{lem:D}.
\end{proof}

\begin{corollary}\label{cor:where}
Let $R_+=\{\lambda_1,\lambda_2>0>\lambda_3\}$ (beyond the side $p_1p_2$, inside the angle at $p_3$) and
$R_-=\{\lambda_1,\lambda_2<0<\lambda_3\}$ (the vertical angle at $p_3$). Every critical point lies in
$R_+\cup R_-$, and those in $R_-$ are non-degenerate saddles. Minima, and degenerate critical points, lie in
\[
  R_+^{\mathrm{out}}=R_+\cap\{P_x<0\}.
\]
\end{corollary}

\begin{proof}
By Lemma~\ref{lem:type}(2), the signs of $\lambda(x)$ are $\pm(+,+,-)$.

In $R_-$, $\operatorname{sign}\sin\sigma_i=\operatorname{sign}\lambda_i$, so $\prod\sin\sigma_i>0$. Moreover $P_x<0$:
$p_3$ is a convex combination of $x,p_1,p_2$ with positive weights, so if $x$ were in the closed circumdisc, $p_3$
would lie in the open disc. By (D), $D<0<c$.

In $R_+$, $\prod\sin\sigma_i<0$, so $D\ge c>0$ forces $P_x<0$.
\end{proof}

\section{Angle coordinates and the gradient structure}\label{sec:grad}

For $x\in R_+$ the segment $p_3x$ crosses $p_1p_2$. So $\sigma_1=t_1\in(0,\pi)$, $\sigma_2=t_2\in(0,\pi)$ and
$\sigma_3=-s$, where $s=t_1+t_2$. On $R_+^{\mathrm{out}}$, (P) gives $\sin(t_i-A_i)<0$, so
\[
  t\in Q:=(0,A_1)\times(0,A_2).
\]

\begin{lemma}\label{lem:inj2}
The map $x\mapsto t=(t_1,t_2)$ is injective on $R_+^{\mathrm{out}}$.
\end{lemma}

\begin{proof}
The point $x$ is the second intersection, other than $p_3$, of the circle through $p_2,p_3$ determined by $t_1$ and the
circle through $p_3,p_1$ determined by $t_2$. These circles are distinct: otherwise their common circle would be the
circumcircle, which does not contain $x$.
\end{proof}

By (P), $r_i=r_1r_2r_3|\sin(\sigma_i-A_i)|/(|P_x|\sin A_i)$. Hence
\[
  \lambda_ir_i^\beta=C(x)\,\sin\sigma_i\,|\sin(\sigma_i-A_i)|^b/\sin^bA_i,\qquad C(x)>0 .
\]
Put
\[
  \varphi_i(t)=\log\sin t+b\log\sin(A_i-t)\ (i=1,2),\qquad \varphi_3(s)=\log\sin s+b\log\sin(s+A_3),
\]
and
\[
  H(t)=\Phi_1(t_1)+\Phi_2(t_2)-\Phi_3(t_1+t_2),\qquad \Phi_i'=\varphi_i .
\]
Then $x\in R_+^{\mathrm{out}}$ is critical for $q$ if and only if $\nabla H(t(x))=\psi_q$, where
\[
  \psi_q=\bigl(\log\tfrac{q_1}{|q_3|}+b\log\tfrac{\sin A_1}{\sin A_3},\ \log\tfrac{q_2}{|q_3|}+b\log\tfrac{\sin A_2}{\sin A_3}\bigr).
\]

\section{The region of minima outside the triangle}

Let $U_i(t)=\cot t+\cot(a_i-t)$ with $a_1=A_1$, $a_2=A_2$, $a_3=A_1+A_2$, and evaluate $U_3$ at $s$. In the
coordinates $t$ we have $W=(-U_1,-U_2,U_3)$, so by (E)
\begin{equation}\label{eq:E}
  E:=U_1U_2-U_3(U_1+U_2)=1/D .
\end{equation}
Thus minima and degenerate points correspond to points of $\overline{M'}$, where $M'=\{t\in Q:E(t)<K\}$.

\begin{lemma}\label{lem:sign}
\begin{enumerate}
\item $E\ge4$, with equality exactly at $t^\circ=(A_1/2,A_2/2)$, which corresponds to the excentre opposite $p_3$.
\item With $d=(\varphi_1',\varphi_2',-\varphi_3')$ we have $e_2(d)=(b+1)^2-bE$. So $\Hess H$ is definite on $M'$; if
  $M'$ is connected, it is negative definite there.
\end{enumerate}
\end{lemma}

\begin{proof}
Put $\tilde x=(\cot t_1,\cot t_2,-\cot s)$ and $\tilde y=(\cot(A_1-t_1),\cot(A_2-t_2),-\cot(A_1+A_2-s))$. Both satisfy
$e_2=1$, by the cotangent addition formula. Also $\tilde x+\tilde y=(U_1,U_2,-U_3)$ and $d=\tilde x-b\tilde y$.

Polarising $e_2$ gives $E=2+2\langle\tilde x,\tilde y\rangle$ and $e_2(d)=1+b^2-b(E-2)$, where $\langle\cdot,\cdot\rangle$
is the Lorentzian polarisation of $e_2$ (Remark~\ref{rem:hyp}). Moreover $\tilde x$ and $\tilde y$ are the cotangent
vectors of the genuine triangles with angles $(t_1,t_2,\pi-s)$ and $(A_1-t_1,A_2-t_2,A_3+s)$. So both lie on the sheet
$\{e_2=1,\ \sum>0\}$ of the hyperboloid, hence $\langle\tilde x,\tilde y\rangle\ge1$, with equality iff
$\tilde x=\tilde y$, i.e.\ iff $t=t^\circ$. (Only $E(t^\circ)=4<K$ is used below.)

At $t^\circ$ we have $d=\bigl((1-b)\cot\tfrac{A_1}2,(1-b)\cot\tfrac{A_2}2,(b-1)\cot\tfrac{A_1+A_2}2\bigr)$, which has
two negative entries. With $e_2(d)>0$, this means $\Hess H=\operatorname{diag}(d)|_{\sum v=0}$ is negative definite.
\end{proof}

\begin{lemma}[index count]\label{lem:index2}
If $V|_\Pi$ is Morse, then $\#\{\text{saddles}\}=\#\{\text{extrema}\}+2$ when $q_1+q_2+q_3\ne0$, and
$\#\{\text{extrema}\}+1$ otherwise.
\end{lemma}

\begin{proof}
We apply Poincar\'e--Hopf to $\nabla V$ on a large disc minus small discs around the charges.
\begin{itemize}
\item On each small circle, $\nabla V$ is radial, so it has degree $1$ whatever the sign of the charge.
\item On a large circle the degree is $1$ when the total charge is nonzero.
\item When the total charge is zero, the dipole moment $m=\sum q_ip_i$ is nonzero, and the field is asymptotic to
  $\nabla(\langle m,x\rangle|x|^{-\gamma})$, $\gamma=2\alpha+2>2$. On circles this field is proportional to
  $(1-\tfrac\gamma2)-\tfrac\gamma2e^{2i\phi}$, which has degree $2$.
\end{itemize}
Extrema have index $+1$ and saddles index $-1$.
\end{proof}

\section{An explicit description of the region}

Put $\Lambda_i(t)=\log\bigl(\sin t/\sin(a_i-t)\bigr)$. Then $\Lambda_i'=U_i$,
$U_i=2(\cosh\Lambda_i+\cos a_i)/\sin a_i$, and $t=\arg(1+e^{\Lambda+ia})$.

\begin{lemma}[product representation]\label{lem:product}
With $\Delta:=\Lambda_1(t_1)-\Lambda_2(t_2)$ (from here on, $\Delta$ denotes this difference, not the area),
\[
  E=\frac{2U_3(s)}{\sin A_3}\,(\cosh\Delta+\cos A_3)=\frac{2\,(\cosh\Delta+\cos A_3)}{\sin s\,\sin(s+A_3)} .
\]
Similarly
\[
  E=\frac{2\,\bigl(\cosh(\Lambda_2(t_2)-\Lambda_3(s))-\cos A_1\bigr)}{\sin t_1\,\sin(A_1-t_1)},\qquad
  \Lambda_3(s)=\log\frac{\sin s}{\sin(A_1+A_2-s)} .
\]
\end{lemma}

\begin{proof}
Geometrically, $E=2+2\cosh\rho$, where $\rho$ is the hyperbolic distance between the shapes $\tilde x,\tilde y$. The
formula is the distance formula of the strip model of the hyperbolic plane, anchored at vertex $3$ (resp.\ $1$).

For a direct proof, write every sine and cosine through $e^{it_1}$, $e^{it_2}$, $e^{iA_1}$, $e^{iA_2}$. Both identities
then become identities between rational functions, and they are verified exactly in
\texttt{code/symbolic\_product.py}.
\end{proof}

Since $1/(\sin s\sin(s+A_3))=U_3/\sin A_3$, Lemma~\ref{lem:product} gives
\begin{equation}\label{eq:region}
  E<K\iff \cosh\Delta<\beta(s):=\tfrac12\bigl[(F-1)\cos A_3-(F+1)\cos(2s+A_3)\bigr].
\end{equation}

\begin{lemma}\label{lem:disc}
\begin{enumerate}
\item The map $t\mapsto(s,\Delta)$ is a diffeomorphism of $Q$ onto $(0,A_1+A_2)\times\R$.
\item Write $2s+A_3=\pi+\psi$. Then $\beta\ge1$ exactly when $|\psi|\le\psi_{\max}$, where
  $\cos\psi_{\max}=(2-(F-1)\cos A_3)/(F+1)$, and $\psi_{\max}<A_1+A_2$.
\item Consequently $M'=\{|\Delta|<\arccosh\beta(s)\}$ is a topological disc, and $\partial M'=\{\cosh\Delta=\beta(s)\}$
  is a smooth simple closed curve.
\end{enumerate}
\end{lemma}

\begin{proof}
(1) For fixed $s$, $\Delta$ increases strictly with $t_1$ along $t_1+t_2=s$, from $-\infty$ to $+\infty$. Its derivative
is $U_1+U_2>0$.

(2) The bound $\psi_{\max}<A_1+A_2=\pi-A_3$ is equivalent to $\cos A_3>-1$.

(3) The gradient $(-\beta',\sinh\Delta)$ does not vanish on the curve: if $\Delta=0$, then $\beta=1<\max\beta$, so
$\beta'\ne0$.
\end{proof}

Given $(s,\Delta)$, the point $t$ is explicit: $\zeta=e^{\Lambda_1}$ is the positive root of
\begin{equation}\label{eq:zeta}
  e^{-\Delta}\sin(A_1+A_2-s)\,\zeta^2+\bigl[\sin(A_1-s)+e^{-\Delta}\sin(A_2-s)\bigr]\zeta-\sin s=0 ,
\end{equation}
which expresses $\arg(1+\zeta e^{iA_1})+\arg(1+\zeta e^{-\Delta}e^{iA_2})=s$.

\section{The curvature criterion}

\begin{lemma}\label{lem:curv}
In linear coordinates $X$, let $\hat G(X)=\beta(X_1+X_2)-\cosh\bigl(\tilde\Lambda_1(X_1)-\tilde\Lambda_2(X_2)\bigr)$.
Write $\tilde U_i=\tilde\Lambda_i'$ (of any sign), $S=\sinh\Delta$, $C=\cosh\Delta$ and $P=-\beta''$. On $\{\hat G=0\}$,
the vector $\tau=(\beta'+S\tilde U_2,\,S\tilde U_1-\beta')\ne0$ is tangent, and
\[
  -\Hess\hat G(\tau,\tau)=J:=(\tilde U_1+\tilde U_2)^2(PS^2+C\beta'^2)
    +S\bigl[\tilde U_1'(\beta'+S\tilde U_2)^2-\tilde U_2'(S\tilde U_1-\beta')^2\bigr].
\]
Hence the boundary of $\{\hat G>0\}$ has positive curvature at a point if and only if $J>0$ there. In particular,
$\{\hat G>0\}$ is strictly convex near such a point.
\end{lemma}

\begin{proof}
We have $\nabla\Delta=(\tilde U_1,-\tilde U_2)$ and $\Hess\Delta=\operatorname{diag}(\tilde U_1',-\tilde U_2')$, so
\[
  \Hess\hat G=\beta''\mathbf 1\mathbf 1^{\mathsf T}-C\nabla\Delta\nabla\Delta^{\mathsf T}-S\Hess\Delta .
\]
The vector $\tau$ is orthogonal to $\nabla\hat G=\beta'\mathbf 1-S\nabla\Delta$. Moreover
$\tau_1+\tau_2=S(\tilde U_1+\tilde U_2)$ and $\tilde U_1\tau_1-\tilde U_2\tau_2=\beta'(\tilde U_1+\tilde U_2)$.
Substituting gives $J$.
\end{proof}

We use Lemma~\ref{lem:curv} in two coordinate systems.

\emph{Vertex~3.} Take $X=t$, $\tilde U_i=U_i>0$, $\tilde\Lambda_i=\Lambda_i$. Dividing by $(U_1U_2)^2$,
\begin{equation}\label{eq:I}
  I=(c_1+c_2)^2(PS^2+C\beta'^2)+S\bigl[a_1(\beta'c_2+S)^2-a_2(S-\beta'c_1)^2\bigr],
\end{equation}
where $c_i=1/U_i=\sin A_i/(2h_i)$, $a_i=U_i'/U_i^2=\sinh\Lambda_i/h_i$ and $h_i=\cosh\Lambda_i+\cos A_i$. The expression
$I$ is invariant under exchanging the indices $1,2$ together with $S\mapsto-S$. So the arc $\Delta<0$ of a triangle is
the arc $\Delta>0$ of the mirrored triangle.

\emph{Vertex~1.} Take $X=(s,-t_2)$, $\tilde\Lambda_1=\Lambda_3$ and $\tilde\Lambda_2(X_2)=\Lambda_2(-X_2)$, so that
$\tilde U_1=U_3$ and $\tilde U_2=-U_2$. Use the representation $\cosh(\Lambda_3-\Lambda_2)<\beta_1(t_1)$, where
$\beta_1=\tfrac12[(F+1)\cos\psi_1-(F-1)\cos A_1]$ and $\psi_1=2t_1-A_1$. Dividing by $U_2^2$,
\begin{equation}\label{eq:J1}
  J_1/U_2^2=(1-U_3c_2)^2(PS^2+C\beta_1'^2)+S\bigl[U_3'(\beta_1'c_2-S)^2-a_2(SU_3-\beta_1')^2\bigr],
\end{equation}
with $U_3'=-U_3(\cot s+\cot(s+A_3))$.

\begin{remark}
These formulas were also tested against the curvature of $\{E=K\}$ computed directly from $E$. At $\alpha=10$, where
$M'$ is not always convex, the signs agree at every one of several thousand boundary points.
\end{remark}

\section{The computer-assisted step}\label{sec:cert}

\begin{proposition}\label{prop:M}
For every $\alpha\in[\tfrac14,2]$ and every triangle, $\partial M'$ has strictly positive curvature. Hence
$\overline{M'}$ is strictly convex.
\end{proposition}

\begin{proof}[Description of the verification]
Parametrise the boundary by $\rho\in[-1,1]$, with $\psi=\rho\psi_{\max}$ and the arc sign of $\Delta$. Then $S$,
$C=\beta$, $\beta'=-(F+1)\sin\psi$ and $P=2(F+1)\cos\psi$ are explicit. Moreover
$\beta-1=(F+1)\sin\frac{\psi_{\max}(1+\rho)}2\sin\frac{\psi_{\max}(1-\rho)}2$, and $\zeta$ is given by \eqref{eq:zeta}.
We cover the space of triangles by three charts.

\emph{Main chart.} $A_1+A_2\in[e_0,\pi]$ and $A_1,A_2\le\pi-r_0$; we verify $I>0$ on the arc $\Delta>0$. The other arc
follows by the mirror symmetry.

\emph{Band chart.} $A_1+A_2=e\in[0,e_0]$, the small triangles. Here $I=O(e^4)$, and we verify $I/e^4>0$ in scaled
variables that are regular at $e=0$:
\begin{itemize}
\item $\psi_{\max}=2\arcsin\bigl(q\sin\frac e2\bigr)$ with $q=\sqrt{(F-1)/(F+1)}$;
\item $S,\beta',c_i$ and $\Delta$ are each $e$ times an analytic function, written with $\frac{\sin x}x$,
  $\frac{\arcsin x}x$ and $\frac{\operatorname{arcsinh}x}x$;
\item the difference $(a_1-a_2)/e$ is enclosed by the mean value theorem.
\end{itemize}

\emph{Corner chart.} $A_1=\pi-r$ with $r\le r_0$, $A_2=\theta r$ and $A_3=(1-\theta)r$. In the vertex-$3$ form some
quantities blow up here, so we verify \eqref{eq:J1} on both arcs, where every quantity stays bounded. The boundary point
is explicit:
\begin{itemize}
\item $\zeta=e^{\Lambda_2}$ is the positive root of
  $e^{\Delta}\sin(A_1-t_1)\zeta^2+[e^{\Delta}\sin(A_1+A_2-t_1)-\sin(A_2+t_1)]\zeta-\sin t_1=0$;
\item $t_2=\arg(1+\zeta e^{iA_2})$, which equals $\arctan\bigl(\zeta\sin A_2/(1+\zeta\cos A_2)\bigr)$ here, since
  $A_2\le r_0<\pi/2$.
\end{itemize}
The corner $A_2\ge\pi-r_0$ is the mirror image of this chart.

The computation uses interval arithmetic in IEEE binary64:
\begin{itemize}
\item every operation is rounded outward by two units in the last place, and every library transcendental by four,
  assuming the math library is accurate to two units;
\item each box of parameters receives an enclosure of the quantity to be verified, and boxes whose lower bound is not
  positive are bisected.
\end{itemize}
The main chart takes $A_1+A_2$ up to the double just above $\pi$, so that triangles with arbitrarily small $A_3>0$ are
covered.

For Coulomb, with $(e_0,r_0)=(0.2,0.05)$, all boxes are certified: $127{,}403$ (main), $2{,}048$ (band) and
$2\times1{,}024$ (corner), in a few seconds. The same holds for $(e_0,r_0)=(0.1,0.1)$ and $(0.4,0.02)$.

For Coulomb, every certified box was also re-verified with the multiple-precision interval arithmetic of
\texttt{mpmath.iv} (80-bit precision), using an independent implementation of the formulas. That run does not rely on
the accuracy assumption for the math library.

\emph{Other exponents.} The verified quantities depend on $\alpha$ only through $F$. Since
$F=(b+1)^2/(2b)-1$ increases from $F(\tfrac14)=\tfrac{13}{12}$ to $F(2)=\tfrac{13}5$, we add $F$ as a fourth interval
variable. The steps are as follows.
\begin{itemize}
\item The interval $[\tfrac14,2]$ is cut into $163$ slices of $\alpha$, of widths from $0.05$ down to $2^{-9}$, narrower
  for larger $\alpha$.
\item On each slice, the three charts are verified with $(e_0,r_0)=(0.4,0.05)$, for all $F$ in a floating-point box
  around the values of $F$ on the slice. In total there are about $1.3\times10^9$ boxes, and the computation takes
  about four hours on one core.
\item A separate run covers the strip of the main chart with $A_1+A_2$ between $\pi-10^{-6}$ and the double just above
  $\pi$, for all $F\in[\tfrac{13}{12},\tfrac{13}5]$.
\item Finally, we check that the certified $F$-boxes cover $[\tfrac{13}{12},\tfrac{13}5]$. The two ends are compared
  with these rationals exactly, so no error analysis of $F$ in floating point is needed.
\end{itemize}
The enclosures with interval $F$ were also tested against an independent evaluation at $40$--$60$ digits, which locates
the boundary point by bisection. At about $4{,}000$ random points, every value was contained in its enclosure.

Unlike the Coulomb case, the computation for other exponents was not repeated in multiple-precision arithmetic. It
relies on the accuracy assumption for the math library. The code is listed in Section~\ref{sec:code}.
\end{proof}

\section{Proof of Theorems~\ref{thm:mixed} and~\ref{thm:alpha}}\label{sec:proof2}

\emph{(a) Non-collinear charges.} By Lemma~\ref{lem:type}(1) and Corollary~\ref{cor:where}, every critical point lies
in $\Pi$, those in $R_-$ are saddles, and the minima and degenerate points lie in $R_+^{\mathrm{out}}$.

By Lemma~\ref{lem:inj2} and Section~\ref{sec:grad}, these correspond injectively to points $t\in\overline{M'}$ with
$\nabla H(t)=\psi_q$.

Suppose there were two such points, $t\ne t'$. By Proposition~\ref{prop:M} the open segment between them lies in $M'$,
and by Lemma~\ref{lem:sign} $\Hess H\prec0$ there. Hence
\[
  \langle\nabla H(t)-\nabla H(t'),t-t'\rangle=\int_0^1(t-t')^{\mathsf T}\Hess H\,(t-t')\,d\lambda<0,
\]
which contradicts $\nabla H(t)=\nabla H(t')$. So there is at most one minimum or degenerate critical point.

For the saddles, Step~3 of Section~\ref{sec:proof1} applies, with $\log|Q_i/Q_3|$ in place of $\log(Q_i/Q_3)$,
with $\Psi$ defined on $R_+\cup R_-$, and with $\kappa$ of either sign. Degenerate critical points are those at which
$D\Psi$ is singular, so Morse charges are dense. If $V$ had four saddles, they would persist for a nearby Morse charge
vector. That contradicts Lemma~\ref{lem:index2} and the bound on extrema. Altogether there are at most four critical
points.

\emph{(b) Collinear charges.} Off the line, $\nabla V=0$ forces $\sum w_i=0=\sum w_it_i$, where $t_i$ are the positions
on the line. By rotational symmetry about the line, such equilibria form circles in $\R^3$, and $(n-2)$-spheres in
$\R^n$.

A non-degenerate equilibrium on the line has Hessian $\operatorname{diag}(f'',h,h)$, so it is also non-degenerate for
$V|_\Pi$, for any plane $\Pi$ through the line. Moving $p_3$ slightly off the line inside $\Pi$, these equilibria
persist as distinct critical points of a non-collinear configuration (implicit function theorem). By (a) there are at
most four. \qed

\section{Remarks on Part~II}\label{sec:rem2}

\begin{remark}[large exponents]\label{rem:large}
For large $\alpha$, $M'$ is not convex for some triangles. For example, at $\alpha=10$ the boundary of $M'$ has negative
curvature for the triangle with angles $(10.6^\circ,29.3^\circ,140.1^\circ)$.

Numerically, the scale-free quantity $I/(A_1+A_2)^4$ is positive for every triangle when $\alpha\le5.6$. It becomes
negative for some nearly degenerate triangles, with one angle close to $\pi$, when $\alpha\ge5.8$. These values were
confirmed at $40$--$50$ digits (\texttt{convexity\_threshold.py}). So between $\alpha=2$ and about $5.7$, the range of
Theorem~\ref{thm:alpha} is limited by computing time rather than by a failure of convexity. Near $\alpha\approx5.75$ the
margin tends to $0$.

At the other end, as $\alpha\to0$ the level $K$ decreases to $4=\min E$, so $M'$ shrinks to the point $t^\circ$. The
charts would have to be rescaled around $t^\circ$ there.

Even beyond this threshold, the charge map remains injective on $M'$ numerically: the image of $\partial M'$ winds at
most once around any point. So Maxwell's bound for three charges seems to hold for all $\alpha$, but a different
argument would be needed to prove it.
\end{remark}

\begin{remark}
Numerically, three collinear charges of different signs have $0$ or $2$ equilibria on the line.

Part (b) uses $n\ge3$. In $\R^2$ the off-axis equilibria of collinear charges are isolated mirror pairs, and they are
not covered by (b).

Part (b) counts only non-degenerate equilibria, and isolated degenerate ones do occur. For charges $-2,24,-64$ at the
points $0,1,3$ of a line in $\R^3$, the point $-1$ of the line is a critical point with vanishing Hessian. No circle of
critical points comes near it, so it is an isolated equilibrium.
\end{remark}

\section{Code}\label{sec:code}

The code is available, in the folder \texttt{code}, at
\begin{center}
\footnotesize\url{https://github.com/anshM123/Maxwell-s-Problem-on-the-Number-of-Equilibria-of-Point-Charges}
\end{center}
The numerical checks of Part~I are \texttt{key\_identities.py}, \texttt{separable.py},
\texttt{predict\_test.py} and \texttt{scan.py}. The verification for Part~II consists of:
\begin{itemize}
\item \texttt{ia.py}: interval arithmetic;
\item \texttt{mixed\_cert.py}: the main chart;
\item \texttt{mixed\_charts.py}: the band and corner charts;
\item \texttt{mixed\_cert\_all.py}: the driver;
\item \texttt{mixed\_certF.py}: other exponents, with $F$ as an interval variable;
\item \texttt{mixed\_alpha\_sweep.py}: the sweep over $\alpha$;
\item \texttt{mixed\_certF\_sliver.py}: the strip next to $A_1+A_2=\pi$, for all $F$;
\item \texttt{sweep\_coverage.py}: checks that the certified $F$-boxes cover $[\tfrac{13}{12},\tfrac{13}5]$;
\item \texttt{check\_certF\_containment.py}: a test of the enclosures for interval $F$ against an independent
  high-precision evaluation;
\item \texttt{convexity\_threshold.py}: the numerics of Remark~\ref{rem:large};
\item \texttt{symbolic\_mixed.py}: the identities of Lemma~\ref{lem:signed};
\item \texttt{symbolic\_product.py}: the identities of Lemma~\ref{lem:product}.
\end{itemize}

\begin{thebibliography}{99}

\bibitem{ABK} P. Arathoon, G. Ball, and M. D. Kvalheim, \emph{The Maxwell conjecture is false}, arXiv:2607.27197 (2026).

\bibitem{EFO} H. Edelsbrunner, C. Fillmore, and G. Oliveira, \emph{Counting equilibria of the electrostatic potential},
Proc. Lond. Math. Soc. (3) \textbf{132} (2026), no.~6, Paper No. e70163.

\bibitem{EP} A. Enciso and D. Peralta-Salas, \emph{The finiteness conjecture for equilibria of electric fields generated
by point charges of one sign}, arXiv:2609.03965 (2026).

\bibitem{GNS} A. Gabrielov, D. Novikov, and B. Shapiro, \emph{Mystery of point charges}, Proc. London Math. Soc. (3)
\textbf{95} (2007), 443--472.

\bibitem{GNNS} A. Gabrielov, D. Novikov, T. Novikov, and B. Shapiro, \emph{From 12 to 6: sharpening the three-charge
bound in Maxwell's problem}, arXiv:2607.28785 (2026).

\bibitem{Johnson} R. A. Johnson, \emph{Advanced Euclidean Geometry}, Dover, New York, 1960.

\bibitem{Max} J. C. Maxwell, \emph{A Treatise on Electricity and Magnetism}, Vol.~1, Clarendon Press, Oxford, 1873.

\bibitem{Oli} G. Oliveira, \emph{Critical points of point charge potentials along lines}, arXiv:2609.00336 (2026).

\bibitem{Tsai} Y.-L. Tsai, \emph{Maxwell's conjecture on three point charges with equal magnitudes}, Phys. D
\textbf{309} (2015), 86--98.

\bibitem{Zol} V. Zolotov, \emph{Upper bounds for the number of isolated critical points via the Thom--Milnor theorem},
Anal. Math. Phys. \textbf{13} (2023), Paper No. 81.

\end{thebibliography}

\end{document}
